% Laget automatisk av elevark/lag_elevark.py fra sim/ballstabel.html. Endringer her blir overskrevet.
\documentclass[11pt]{article}
\usepackage{elevark}
\begin{document}
\elevarkhode{Fysikk 1 · Bevegelsesmengde og støt}{Ballstabel}{https://lektorbo.no/sim/ballstabel.html}{Simuleringen}

Legg en tennisball oppå en basketball og slipp dem samtidig fra 1~m. Basketballen spretter fra gulvet og treffer tennisballen, som fortsatt er på vei ned. Hvor høyt spretter tennisballen?

\section*{Oppgaver}

Prøv gjerne med ekte baller først. Hold tennisballen rett over basketballen, og slipp begge samtidig utendørs eller i en gymsal.

\begin{oppgaveliste}
\oppgave \fase{gjett}{Gjett}Basketball og tennisball slippes fra 1~m. Hvor høyt spretter tennisballen? Lavere enn 1~m, omtrent 1~m, 2–3~m eller enda høyere?
\svar
\oppgave \fase{test}{Test}Kjør simuleringen med ideelle støt. Hvor høyt kommer tennisballen? Hva skjer med basketballen?
\svar
\oppgave \fase{forklar}{Forklar}Se på støtet fra basketballens synspunkt. Rett før støtet går basketballen opp med farten $v$ og tennisballen ned med farten $v$. Hvor fort nærmer tennisballen seg basketballen? Hvor fort går den fra basketballen etter et elastisk støt mot en mye tyngre ball? Hvor fort går den da sett fra gulvet?
\svar
\oppgave \fase{test}{Regn}Bruk resonnementet fra oppgave 3 til å vise at en veldig lett ball på en veldig tung ball får farten $3v$ og spretter 9 ganger så høyt. Sjekk med bordtennisballen i simuleringen. Hvorfor blir det ikke helt 9?
\svar
\oppgave \fase{forklar}{Forklar}Tennisballen får mer kinetisk energi enn den hadde før. Bryter det med energibevaring? Hvor kommer energien fra?
\svar
\oppgave \fase{test}{Undersøk}Fjern avkryssingen for ideelle støt. Hvor mye lavere blir sprettet? Prøv to tennisballer oppå hverandre. Hva skjer nå, og hvorfor?
\svar
\oppgave \fase{utfordring}{Utfordring}Tenk deg en stabel med tre baller der hver ball er mye lettere enn den under. Hvor fort går den øverste ballen, og hvor høyt kan den komme? Tips: Bruk samme resonnement to ganger.
\svar
\end{oppgaveliste}

\end{document}
